\chapter{为什么要写这本书}
\label{ch:overview}

原始手稿提出了一个有用的核心判断：人之所以经常无法启动那些明知有益的行动，并不只是因为他们缺少价值判断，而是因为他们被嵌在某个既有的行为框架里，这个框架本身的回报结构、认知惯性与所处环境都会共同阻碍转变。这个判断仍然是整个项目的中心。本稿改变的，是它的教学目标。

这一版写成的是教材原型，而不再是一篇较短的概念性文章。教材需要的，不只是方向正确的直觉。它还需要有序的定义、稳定的记号、层层推进的章节结构、展开后的案例，以及一套能让读者检验这个框架究竟是否真能解释行为，而不只是听起来有启发性的练习层。

\section{学习目标}

读完本稿后，读者应当能够：
\begin{enumerate}[nosep]
  \item 把行为状态描述成一个结构化对象，而不只是某种情绪；
  \item 区分回报、切换成本与环境强化；
  \item 识别自由意志具有异常高杠杆作用的决策窗口；
  \item 解释为什么计划改变的是后续行动的几何结构，而不只是额外提供一点动力；
  \item 把行为转变与关于意识状态稳定性及状态变化的实证研究联系起来；
  \item 解释为什么框架会从重复数据、耦合与干扰中涌现出来，而不是以静态标签的方式直接出现；
  \item 判断哪些高阶动力系统工具能够澄清这个框架，哪些内容只适合放在附录或前沿讨论中；
  \item 用一套可复用的案例语言，对现实中的失败进行分类；
  \item 在失败发生之前，设计能够提前降低激活障碍的干预方案。
\end{enumerate}

\section{全书架构}

这份手稿围绕九项教学任务展开：
\begin{enumerate}[nosep]
  \item 第~\ref{ch:formal-language} 章与第~\ref{ch:dynamics} 章定义模型的语言。
  \item 第~\ref{ch:planning} 章解释为什么计划在结构上不可或缺。
  \item 第~\ref{ch:consciousness-biology} 章引入意识层的实证材料，并把状态转变与生物学证据连接起来。
  \item 第~\ref{ch:emergence} 章解释持久框架如何在重复互动中形成，而不是来自某一次孤立选择。
  \item 第~\ref{ch:advanced-dynamics} 章把数学语言从轻量级的障碍模型升级为更丰富的状态空间视角。
  \item 第~\ref{ch:cases} 章用多个领域的案例把模型具体化。
  \item 第~\ref{ch:interventions} 章把分析转成可执行的设计模式。
  \item 第~\ref{ch:exercises} 章检验读者是否能够复用这一框架。
  \item 附录保留记号、模板，以及与主体内容严格分离的前沿假说。
\end{enumerate}

\section{范围与边界}

这仍然是一个半形式化项目。它的目的，是在数学化表达与行为解释力之间搭起一座稳固的桥。因此：
\begin{itemize}[nosep]
  \item 记号体系有纪律，但刻意保持轻量；
  \item 这里的命题是解释工具，而不是已经定稿的实证定律；
  \item 实证神经科学、理论系统语言与带有猜想性质的意识物理提案，被放在不同的证据通道里；
  \item 案例库的设计目标是提升迁移能力，而不是替代心理学文献；
  \item 干预模式是为后续扩展准备的生产模板。
\end{itemize}

\section{证据政策}

这个有研究支撑的版本采用四条证据通道：
\begin{enumerate}[nosep]
  \item \textbf{强实证核心：} 适合支撑主线意识论证与状态转变论证的材料。
  \item \textbf{机制构建型理论：} 能够加深解释、但并不冒充直接测量结果的涌现与动力学文献。
  \item \textbf{前沿猜想：} 风险高、并且明确处于实证核心之外的提案。
  \item \textbf{推荐阅读：} 便于读者逐章深入的经典教材与重要综述。
\end{enumerate}

\begin{readingbox}
\textbf{基础}
\begin{itemize}[nosep]
  \item Stanislas Dehaene, \emph{Consciousness and the Brain}.
  \item Christof Koch, \emph{Consciousness}.
\end{itemize}
\textbf{现代研究}
\begin{itemize}[nosep]
  \item 未来章节继续扩展前，应先使用分级文献地图来精确筛选论文。
\end{itemize}
\textbf{前沿}
\begin{itemize}[nosep]
  \item 只有在理解实证核心之后，才适合进入前沿阅读。
\end{itemize}
\end{readingbox}

\begin{summarybox}
这一章改变了整份手稿的功能定位。这个项目不再只是一个单一论点，而是一套可复用的教学系统，拥有明确的学习路径、清晰的边界，以及一种可以持续吸收后续证明、案例与练习而不坍塌成零散笔记的结构。
\end{summarybox}
